\begin{definition}[Carnot group with its homogeneous norm]
A \emph{Carnot norm datum} is a Carnot group datum (\noderef{CarnotGroup}) together with a
map $\|\cdot\|:\mathbb R^n\to\mathbb R$, the homogeneous norm, subject to the conditions
below. Throughout, $q:=2\,r!$, the associated left-invariant distance is
\[
 d(x,y):=\|y^{-1}*x\|=\|(-y)*x\|,
\]
and $B_\rho(x):=\{y\in\mathbb R^n: d(y,x)<\rho\}$.
\begin{enumerate}
\item $\|x\|\ge0$ for all $x$; $\|x\|=0$ if and only if $x=0$; $\|-x\|=\|x\|$; and
  $\|\cdot\|$ is continuous on $\mathbb R^n$.
\item $\|\exp_{\mathrm h}(\zeta)\|=|\zeta|$ for every $\zeta\in\mathbb R^m$.
\item $x\mapsto\|x\|^{q}$ is $C^\infty$ on $\mathbb R^n$.
\item (First inequality of (2.17).) For every $\nu>0$ there is $c>0$ such that
  $\|x*y\|\le c\,(\|x\|+\|y\|)$ whenever $\|x\|<\nu$ and $\|y\|<\nu$.
\item (Second inequality of (2.17).) For every $\nu>0$ there is $c>0$ such that
  $\|x^{-1}*y*x\|\le c\,\|y\|^{1/r}$ whenever $\|x\|<\nu$ and $\|y\|<\nu$.
\item (H\"older equivalence on bounded sets.) For every $\nu>0$ there is $c>0$ such that
  \[
   c^{-1}\,|x-y|\ \le\ d(x,y)\ \le\ c\,|x-y|^{1/r}
   \qquad\text{whenever}\quad |x|\le\nu\ \text{and}\ |y|\le\nu .
  \]
\item For every $x\in\mathbb R^n$ and every $\rho>0$ the ball $B_\rho(x)$ is non-empty and
  connected.
\item $\|\cdot\|$ is proper: for every $C\in\mathbb R$ the sublevel set
  $\{x\in\mathbb R^n:\|x\|\le C\}$ is bounded.
\end{enumerate}

These are the properties of the homogeneous norm (2.15) recorded in Section
Section~2.4. Item (1) is continuity, evenness and non-degeneracy of
$\|x\|=\big(\sum_{j}(\sum_i|x_{i,j}|^2)^{r!/j}\big)^{1/2r!}$. Item (3) is the observation
that $\|x\|^{q}=\|x\|^{2r!}$ is a polynomial, which is what makes the kernel
$\|y*x^{-1}\|^{q}$ of the sup- and inf-convolutions (3.13) smooth; item (2) says
that the horizontal exponential is an isometry from $(\mathbb R^m,|\cdot|)$ onto the
first-layer directions, which is immediate from (2.15) since
$\exp_{\mathrm h}(\zeta)$ has first-layer coordinates $\zeta$ and vanishing higher
coordinates. Items (4) and (5) are the two inequalities (2.17), the second being the
estimate $\|c_x(y)\|\le c_{\mathbb G}\|y\|^{1/r}$ for the inner automorphism
$c_x(y)=x^{-1}*y*x$, whose exponent $1/r$ is sharp. Item (6) is the H\"older equivalence of
$d$ with the Euclidean distance on bounded sets, which also says the two topologies agree.
Item (7) is the statement that the balls $B_\rho(0)$ are connected, being star-shaped with
respect to the dilations, and hence so are all balls by left invariance of $d$; it is
non-empty because $d(x,x)=\|0\|=0<\rho$.

Left invariance $d(z*x,z*y)=d(x,y)$ and the right-translation identity (2.16)
$d(x*h,x)=\|h\|$ are not listed: both follow from the group axioms of
\noderef{CarnotGroup} alone, since $d(z*x,z*y)=\|(z*y)^{-1}*(z*x)\|=\|y^{-1}*x\|$ and
$d(x*h,x)=\|x^{-1}*x*h\|=\|h\|$. The dilations themselves are not carried, because nothing
below uses them other than through items (3), (5) and (7).

Item (8) is the properness of (2.15), and it too is immediate from that formula: if
$\|x\|\le C$ then $\big(\sum_i|x_{i,j}|^2\big)^{r!/j}\le C^{2r!}$ for every layer $j$,
hence $\sum_i|x_{i,j}|^2\le C^{2j}$, so all coordinates of $x$ are bounded in terms of $C$
alone. It is listed as a separate condition because it does not follow from the others: the
quasi-triangle inequality (4) bounds $\|y\|$ in terms of $\|y*x^{-1}\|$ and $\|x\|$, which
says nothing about the Euclidean size of $y$, and the lower H\"older bound of (6) is
available only once both of its arguments are already known to be Euclidean-bounded, so
neither of them yields boundedness of a sublevel set of $\|\cdot\|$. Properness is what makes
a sublevel set of the kernel $y\mapsto\|y*x^{-1}\|^{q}$ of (3.15) relatively compact, hence
the infimum in (3.15) locally an infimum over a compact set of $y$ and therefore a
continuous function of $x$; and that continuity is what makes the reduced domain
$\Omega[\varepsilon]$ open, as the paper asserts without proof.

As with \noderef{CarnotGroup}, every Carnot group in the coordinates of Section
Section~2.4 satisfies all eight conditions, so the proviso of item 3 of the
formalization note holds. For the Euclidean additive group on $\mathbb R^m$ ($n=m$, $r=1$,
$x*y=x+y$, $\|x\|=|x|$) one has $q=2$, $d(x,y)=|x-y|$, and each condition reduces to a
standard fact with $c=1$: (4) is the triangle inequality, (5) is $\|y\|\le\|y\|$ since
$x^{-1}*y*x=y$, (6) is an equality, (7) holds because Euclidean balls are convex, and (8) holds because a
Euclidean ball is bounded.
\end{definition}
